% =====================================================================
%  Abstract
% =====================================================================
\begin{abstract}
\noindent
Quantum key distribution (QKD) lets two parties agree on a secret key whose
security rests on the laws of physics rather than on computational
hardness. Moving it from the laboratory into an operator's network,
however, means managing it like any other network resource: discovering its
topology, monitoring its state, allocating its capacity to services and
coordinating it with classical networks. The ETSI ISG QKD addressed this
problem step by step. GS~QKD~002 identifies the use cases and integration
requirements; GS~QKD~015 defines the model through which an SDN controller
sees and configures the QKD nodes of a domain; GS~QKD~018 defines the
interface between an \emph{SDN orchestrator} and the controller of a QKD
domain, enabling end-to-end provisioning across several domains and
between a QKD network and an optical transport network. This work follows
that chain in a reasoned way: it describes the components of an SD-QKD
network and the structure of a YANG model, analyses GS~QKD~018 in detail
(use cases, information model, topology, inventory and connectivity
modules, multi-domain workflow, security), shows its analogy with the IETF
ACTN framework (RFC~8453) and compares it with three real experiences: SK
Telecom's multi-vendor KOREN testbed, an ONOS-based hierarchical prototype
developed at Politecnico di Milano, and the MadQCI network in Madrid. It
closes with the current limits of the standard and open research
directions.

\medskip
\noindent\textbf{Keywords:} QKD, SDN, orchestration, ETSI GS QKD 015,
ETSI GS QKD 018, YANG, RESTCONF, multi-domain, key management.
\end{abstract}
